\documentclass[11pt,a4paper]{article}

% ============================================
% ENCODING & FONTS
% ============================================
\usepackage[utf8]{inputenc}
\usepackage[T1]{fontenc}
\usepackage{lmodern}

% ============================================
% PAGE LAYOUT
% ============================================
\usepackage[margin=1.618cm, top=2.618cm, bottom=2.618cm]{geometry}

% ============================================
% ESSENTIAL PACKAGES
% ============================================
\usepackage{float}
\usepackage{caption}
\usepackage{setspace}
\usepackage{fancyhdr}
\usepackage{xcolor}
\usepackage{hyperref}
\usepackage{amsmath}
\usepackage{amssymb}
\usepackage{amsthm}
\usepackage{booktabs}
\usepackage{array}
\usepackage{graphicx}
\ifdefined\DeclareUnicodeCharacter
  \DeclareUnicodeCharacter{00B0}{\ensuremath{^\circ}}
\fi

% ============================================
% TITLE FORMATTING
% ============================================
\usepackage{titling}
\pretitle{\LARGE\bfseries}
\posttitle{\vspace{-0.4em}}
\preauthor{}
\postauthor{}
\predate{}
\postdate{}
\setlength{\droptitle}{-2.0em}

% ============================================
% HEADER/FOOTER
% ============================================
\setlength{\headheight}{14pt}
\pagestyle{fancy}
\fancyhf{}
\fancyhead[L]{Spectral Geometry of an Explicit $\mathrm{G}_2$ Metric}
\fancyhead[R]{\thepage}
\renewcommand{\headrulewidth}{0.2pt}

% ============================================
% HYPERREF
% ============================================
\hypersetup{
  colorlinks=true,
  linkcolor=blue,
  citecolor=blue,
  urlcolor=blue,
  pdftitle={Spectral Geometry of an Explicit G2 Metric: Laplacian Spectrum and Harmonic Forms},
  pdfauthor={Brieuc de La Fourniere}
}

% ============================================
% SPACING
% ============================================
\setstretch{1.2}
\setlength{\parskip}{0.4em}
\setlength{\parindent}{0pt}

% ============================================
% CUSTOM COMMANDS
% ============================================
\newcommand{\eps}{\varepsilon}
\newcommand{\CIvvv}{\mathrm{CI}(2,2,2)}
\newcommand{\Pfive}{\mathbb{P}^5}
\newcommand{\RR}{\mathbb{R}}
\newcommand{\Gtwo}{\mathrm{G}_2}
\newcommand{\lamone}{\lambda_1}
\newcommand{\Vol}{\mathrm{Vol}}

\pdfstringdefDisableCommands{%
  \def\CIvvv{CI(2,2,2)}%
  \def\Pfive{P^5}%
  \def\Gtwo{G2}%
}

% ============================================
% DOCUMENT
% ============================================

\title{%
\LARGE\textbf{Spectral Geometry of an Explicit $\Gtwo$ Metric:\\[0.2em]
Laplacian Spectrum and Harmonic Forms}
}
\author{}
\date{}

\begin{document}

\maketitle
\noindent\rule{\textwidth}{0.2pt}

\noindent\textbf{Author}: Brieuc de La Fourni\`ere

\noindent Independent researcher, Beaune, France

\noindent Contact: \href{mailto:brieuc@bddelaf.com}{\texttt{brieuc@bddelaf.com}}

\noindent\textbf{Date}: April 2026

\noindent\textbf{DOI}: \href{https://doi.org/10.5281/zenodo.18920368}{10.5281/zenodo.18920368}

\vfill

\begin{abstract}
We compute the Laplacian spectrum and harmonic forms on the certified torsion-free $\Gtwo$ structure of a companion paper~[A], constructed on a TCS-type neck model $K3 \times T^2 \times [0,1]$. The scalar spectral gap is $\lambda_1 = 0.12461 \pm 0.00016$ (Richardson-extrapolated across six grids, Appendix~\ref{app:convergence}) with Weyl exponent $\alpha = 1.998$ (adiabatic prediction: $2.0$), matching the analytical formula $\lambda_1 = \pi^2/(L^2 g_{ss}) = 6\pi^2/475 \approx 0.12467$ --- derived a posteriori from the NK-certified seam parameters --- to $0.05\%$. Spectral analysis of the Hodge Laplacians realizes the harmonic multiplicities $(b_2, b_3) = (21, 77)$ expected from the conjectural compact extension: 21 near-zero eigenvalues of $\Delta_2$ (gap ratio $14{,}635$) and 77 near-zero eigenvalues of $\Delta_3$. The 21 harmonic 2-forms carry a K3 fiber intersection form of signature $(3, 18)$ --- a rank-21 restriction of the ambient K3 lattice $(3, 19)$ --- with SD/ASD eigenvalue gap $2{,}210$. The spectral democracy theorem proven in~[A] is verified numerically to $10^{-5}$ precision. The complete Kaluza--Klein tower comprises $22{,}671$ levels with Poisson level spacing statistics. All spectral results are properties of the explicit neck-model metric; their interpretation on a closed compact $\Gtwo$ manifold is conditional on the existence of such a manifold with $(b_2, b_3) = (21, 77)$.
\end{abstract}

\noindent\textbf{Keywords}: spectral geometry; $\Gtwo$ holonomy; Hodge Laplacian; Kaluza--Klein spectrum; harmonic forms; intersection form.

\vfill
\noindent\rule{\textwidth}{0.2pt}

\newpage
\tableofcontents

\newpage

% ============================================
\section{Introduction}
% ============================================

The construction of explicit metrics with exceptional holonomy has been a long-standing challenge in Riemannian geometry. While existence theorems for compact $\Gtwo$ manifolds are well-established (Joyce~\cite{Joyce2000}, Kovalev~\cite{Kovalev2003}, CHNP~\cite{CHNP2015}, Nordstr\"om~\cite{Nordstrom2023}), explicit spectral data --- whether on closed compact $\Gtwo$ manifolds or on neck-model geometries adapted to them --- has remained unavailable due to the absence of concrete metric coefficients.

In a companion paper~[A], we certify a torsion-free $\Gtwo$ structure on a TCS-type neck model $K3 \times T^2 \times [0,1]$ via computer-assisted proof, with seam-sector geometry adapted to a putative compact extension of Betti numbers $(b_2, b_3) = (21, 77)$, using interval arithmetic with zero finite differences. The present paper computes the spectral geometry of this certified neck-model metric; the spectral results are properties of that metric, while their interpretation on a closed compact $\Gtwo$ manifold is conditional on the existence of such a manifold with $(b_2, b_3) = (21, 77)$ ([A], §2.2). We compute:

\begin{enumerate}
\item \textbf{Scalar Laplacian spectrum} (§\ref{sec:laplacian}): spectral gap, Weyl law, complete KK tower.
\item \textbf{Harmonic forms} (§\ref{sec:harmonic}): spectral Betti confirmation, intersection form, SD/ASD structure.
\item \textbf{Numerical democracy verification} (§\ref{sec:democracy}): confirmation of the analytical theorem from~[A].
\item \textbf{Physical context} (§\ref{sec:discussion}): relevance for Kaluza--Klein compactifications.
\end{enumerate}

\textbf{Note on related work.} An independent arithmetic derivation of the pair $(b_2, b_3) = (21, 77)$ has recently appeared in Zhou \& Zhou~\cite{ZhouZhouA2026,ZhouZhouE2026}, based on Diophantine constraints and spectral stability screening. Their approach is complementary to the present one.

% ============================================
\section{The Certified Metric (Summary)}
% ============================================

We briefly summarize the metric from~[A]. The full construction, NK certification, and analytical decomposition are presented therein.

The metric on a neck region $\mathrm{K3} \times T^2 \times [0,1]$ decomposes into three adiabatic sectors:
\begin{equation}
g(s, \theta, \chi, y) = g_{\mathrm{seam}}(s) \oplus g_{T^2} \oplus g_{K3}(y)
\end{equation}
with effective eigenvalues:
\begin{itemize}
\item $g_{ss} = 19/6$ (torsion-minimizing rational, to $0.03\%$),
\item $g_{T^2} = 7/6$ (torsion-minimizing rational, to $0.20\%$),
\item $g_{K3} \approx 64/77$ (mean K3 eigenvalue, rational approximation to $0.41\%$; see~[A], §3.8),
\item $\det(g) = 65/32$ (exact, imposed at reconstruction; interval-certified to width $1.4 \times 10^{-12}$),
\item $\gamma = 2\pi\sqrt{6/7} \approx 5.817$ (derived from T² Hodge Laplacian + $H^1(\mathrm{K3})=0$; NK-computed: $5.811$, $0.1\%$ proximity).
\end{itemize}

The Newton--Kantorovich certificate~[A] establishes existence and uniqueness of a nearby torsion-free $\Gtwo$-structure with $h = 1.43 \times 10^{-9}$ (margin $\times 350$ million, zero finite differences) and metric correction $\delta g/g \leq 1.35 \times 10^{-7}$.

For the spectral analysis below, the key property is that the torsion-free metric $g^*$ is adiabatically product-type with certified error $\varepsilon_{\mathrm{ad}} = 2 \times 10^{-3}$ ([A], Proposition~3.3), giving eigenvalue perturbation $|\delta\lambda|/\lambda \leq 0.78\%$ and enabling decomposition of 7D Laplacian operators into families of 1D Sturm--Liouville problems.

\begin{quote}
\emph{Remark (independence from the exact value of $g_{K3}$).} The spectral identifications of this paper --- scalar gap $\lambda_1 = 0.1244$, Weyl exponent $\alpha = 1.998$, spectral Betti confirmations $(b_2 = 21, b_3 = 77)$, intersection form signature $(3, 18)$ --- rely only on the adiabatic product structure of $g^*$ with certified error $\varepsilon_{\mathrm{ad}} = 2 \times 10^{-3}$. They do \emph{not} depend on the exact value of $g_{K3}$, on whether $g_{K3} = 64/77$ is an identity or an approximation, or on the precise K3 eigenvalue deviation pattern. The certified perturbation budget $|\delta\lambda_n|/\lambda_n \leq 0.78\%$ absorbs the measured K3 eigenvalue spread of $1.16\%$ and any residual $O(10^{-3})$ uncertainty. The spectral claims herein are robust to future revisions of the analytical status of $g_{K3}$ in~[A].
\end{quote}

% ============================================
\section{Laplacian Spectrum}\label{sec:laplacian}
% ============================================

\subsection{Adiabatic decomposition}

The adiabatic product structure decomposes the 7D Laplacian into families of 1D Sturm--Liouville operators, one per fiber mode $(m, n, \mu_{K3})$:
\begin{equation}
\Delta_0 f = -\frac{1}{\sqrt{\det g}}\, \partial_s \!\left( \sqrt{\det g}\, g^{ss}\, \partial_s f \right) + V_{(m,n,\mu)} f
\end{equation}
where $V_{(m,n,\mu)} = m^2 g^{\theta\theta} + 2mn\, g^{\theta\chi} + n^2 g^{\chi\chi} + \mu_{K3}$ combines the T² and K3 fiber eigenvalues. This decomposition is exact for the product-type metric and certified to $\varepsilon_{\mathrm{ad}} = 2 \times 10^{-3}$ adiabatic error ([A], Proposition~3.3).

We discretize on a Chebyshev collocation grid with $N = 800$ to $1600$ nodes on $s \in [-2, 3]$ and solve the resulting matrix eigenvalue problem using dense diagonalization (Neumann boundary conditions) and shift-invert Lanczos (Dirichlet).

\subsection{Spectral gap}

The fundamental $(0, 0, 0)$ channel yields the scalar spectral gap:
\begin{equation}\label{eq:lambda1}
\lambda_1 = 0.12461 \pm 0.00016
\end{equation}
Richardson-extrapolated across six grid sizes ($N = 200$ to $1200$) with both boundary condition types yielding consistent values (Appendix~\ref{app:convergence}; raw $N=1200$ value $0.12450$).

\begin{table}[ht]
\centering
\begin{tabular}{ccc}
\toprule
$n$ & $\lambda_n$ & $\Delta\lambda/\lambda$ (Richardson) \\
\midrule
0 & $3.5 \times 10^{-13}$ & --- (zero mode) \\
1 & $0.12450$ & $0.08\%$ \\
2 & $0.4976$ & $0.04\%$ \\
3 & $1.1196$ & $0.02\%$ \\
4 & $1.9904$ & $0.01\%$ \\
5 & $3.1099$ & $<0.01\%$ \\
6 & $4.4782$ & $<0.01\%$ \\
7 & $6.0952$ & $<0.01\%$ \\
8 & $7.9609$ & $<0.01\%$ \\
9 & $10.075$ & $<0.01\%$ \\
\bottomrule
\end{tabular}
\caption{First 10 eigenvalues of the $(0,0,0)$ seam channel (Neumann BC, $N=1200$ raw values; the $\Delta\lambda/\lambda$ column gives the Richardson convergence rate, not the deviation from the analytical formula). The Richardson-extrapolated value of $\lambda_1$ is $0.12461 \pm 0.00016$ (eq.~(\ref{eq:lambda1}) and Appendix~\ref{app:convergence}); this is the value used elsewhere in the paper.}
\label{tab:evals}
\end{table}

The near-quadratic growth $\lambda_n \approx 0.125\, n^2$ is characteristic of a 1D Sturm--Liouville problem; the eigenvalue spacing implies an effective spectral length $L_{\mathrm{eff}} \approx \pi/\sqrt{0.125} \approx 8.89$ (distinct from the domain extent $L = 5$ used in the analytical formula below).

\textbf{Analytical formula.} In terms of the NK-certified model parameters, the spectral gap satisfies
\begin{equation}
\lambda_1 = \frac{\pi^2}{L^2 \cdot g_{ss}} = \frac{6\pi^2}{475} \approx 0.12467
\end{equation}
where $L = 5$ is the seam domain extent ($s \in [-2, 3]$, 5 units) and $g_{ss} = 19/6$ is the seam metric component (torsion minimizer, [A] Proposition~3.2), giving $L^2 \cdot g_{ss} = 25 \cdot 19/6 = 475/6$. This is an analytical expression for $\lambda_1$ in terms of NK-certified model parameters, not an independent topological prediction; $L$ and $g_{ss}$ are both outputs of the certification. The match to the Richardson-extrapolated value $0.12461 \pm 0.00016$ (Appendix~\ref{app:convergence}) is $0.05\%$.

\subsection{Weyl law}

The eigenvalue counting function $N(\lambda) = \#\{n : \lambda_n \leq \lambda\}$ satisfies the Weyl asymptotic $N(\lambda) \sim C \lambda^{\alpha/2}$. A least-squares fit to the $(0,0,0)$ channel yields:
\begin{equation}
\alpha = 1.998 \pm 0.002, \qquad C = 0.125
\end{equation}
in agreement with the predicted exponent $\alpha = 2$ for the seam sector. Note that $\alpha = 2$ is the expected Weyl exponent for any 1D Sturm--Liouville problem and reflects the dimension of the seam channel, not a $\Gtwo$-specific feature. The non-trivial spectral verification is the full 7D exponent below.

For the full unified Kaluza--Klein tower (all T² and K3 channels), the effective 7D Weyl exponent is $\alpha_{7D} = 3.460 \pm 0.040$, consistent with $\alpha = 7/2 = 3.5$ (deviation $1.1\%$).

\subsection{Kaluza--Klein tower and level spacing}

The complete KK tower below $\lambda = 20$ contains $22{,}671$ distinct energy levels (comparable in scope to KK spectra computed on weak-$\Gtwo$ examples such as squashed $S^7$~\cite{Nilsson2024}; for a recent comprehensive review of KK reductions in supergravity see~\cite{DuffNilssonPope2025}). The three-scale hierarchy is:
\begin{table}[ht]
\centering
\begin{tabular}{lcc}
\toprule
Sector & First eigenvalue & Scale \\
\midrule
Seam & $0.125$ & $0.353$ \\
$T^2$ & $0.855$ & $0.925$ \\
K3 & $1.208$ & $1.099$ \\
\bottomrule
\end{tabular}
\caption{Sector hierarchy in the KK tower.}
\label{tab:sectors}
\end{table}

The T² channels obey a quadratic potential $V_{(m,n)} = 0.855(m^2 + n^2)$ with isotropy to $3 \times 10^{-7}$. K3 adiabatic additivity holds to $0.003$--$0.023\%$.

The level spacing distribution follows Poisson statistics ($\chi^2 = 1.36$) rather than GOE ($\chi^2 = 28.0$), consistent with an integrable adiabatic system.

% ============================================
\section{Harmonic Forms and Intersection Theory}\label{sec:harmonic}
% ============================================

\subsection{Harmonic 2-forms}

We construct 21 harmonic 2-forms $\omega_I$ on $M$ by lifting K3 harmonic $(1,1)$-forms through the adiabatic decomposition. Each $\omega_I$ has a radial profile $f_I(s)$ satisfying $\Delta_2 \omega_I = 0$, computed via variational optimization with a neural network ansatz~\cite{LarforsLukasRuehle2022,cymyc2024} (3-layer MLP minimizing $\|\Delta_2(f_I\, d\theta \wedge \alpha_I)\|^2_{L^2}$ on the Chebyshev collocation grid, convergence criterion $\max|\mathrm{residual}| < 10^{-7}$) followed by Gram--Schmidt orthonormalization.

The 21 forms decompose into two families:
\begin{itemize}
\item \textbf{$N_1$-type} (11 forms): profiles interpolating from 1 to 0 along the neck.
\item \textbf{$N_2$-type} (10 forms): profiles interpolating from 0 to 1.
\end{itemize}

This decomposition reflects the two-block structure of the neck region, matching the TCS formula $b_2 = b_2(M_1) + b_2(M_2) = 11 + 10$. Three independent decompositions are consistent: $7 + 14$ ($\Gtwo$ representations $\Lambda^2_7 + \Lambda^2_{14}$), $3 + 18$ (SD + ASD from K3 hyperk\"ahler triple), and $11 + 10$ (TCS building blocks). The $L^2$ Gram matrix has condition number $1.047$ with maximum off-diagonal entry $0.012$.

\textbf{Spectral confirmation.} $\Delta_2$ has 21 eigenvalues below $10^{-8}$. The 22nd eigenvalue is $\lambda_{22} \approx 0.12$, giving a gap ratio:
\begin{equation}
\frac{\lambda_{22}}{\max(\lambda_1, \ldots, \lambda_{21})} \approx 14{,}635.
\end{equation}
This four-orders-of-magnitude gap confirms $b_2 = 21$.

\subsection{Intersection structure}

The $L^2$ Gram matrix $G_{IJ} = \int_M \omega_I \wedge \star\omega_J$ is positive definite (condition number $1.047$, §4.1). The topological intersection structure is accessed via the \textbf{K3 fiber pairing} $I_{IJ} = \int_{K3} \alpha_I \wedge \alpha_J$, where $\alpha_I = \omega_I|_{K3}$ are the restrictions to the K3 fiber --- this is the cup-product intersection form on $H^2(K3)$, which takes both positive and negative values.

The ambient K3 lattice $H^2(K3) \cong 3U \oplus 2(-E_8)$ has signature $(3, 19)$ of rank 22. Of the 22 K3 classes, 21 extend to global harmonic forms on $M$ (confirmed by the 21 near-zero $\Delta_2$ eigenvalues, §4.1); the remaining class does not survive TCS boundary matching and corresponds to the first non-zero $\Delta_2$ eigenvalue $\lambda_{22} \approx 0.12$. The K3 fiber pairing restricted to the \textbf{21 global harmonic 2-forms} has signature $\mathbf{(3, 18)}$:

\begin{table}[ht]
\centering
\begin{tabular}{lcl}
\toprule
Type & Count & $I_{IJ}$ eigenvalue range \\
\midrule
Self-dual (SD) & 3 & $4.863,\; 5.499,\; 7.795$ \\
Anti-self-dual (ASD) & 18 & $-0.00423$ to $-0.00219$ \\
\bottomrule
\end{tabular}
\caption{K3 fiber pairing $I_{IJ}$ eigenvalue structure restricted to the 21 global harmonic 2-forms ($3 + 18 = 21 = b_2$). Mean SD / mean ASD ratio: $2{,}210$.}
\label{tab:intersection}
\end{table}

The ratio of mean SD to mean ASD absolute eigenvalue is approximately $\mathbf{2{,}210}$. The cross-block Frobenius norm $\|I_{\mathrm{SD,ASD}}\|_F = 0.018$ confirms near-exact block diagonality.

This extreme gap originates from the $\mathrm{SU}(2)$ holonomy of the K3 fiber: SD 2-forms are aligned with the hyperk\"ahler triple $(\omega_I, \omega_J, \omega_K)$, while ASD forms lie in the 18-dimensional complement of the rank-21 restricted lattice.

\subsection{Harmonic 3-forms}

The 77 harmonic 3-forms decompose as:
\begin{itemize}
\item \textbf{35 constant-type}: $\Lambda^3$-sections with constant profiles.
\item \textbf{21 $d\theta$-fiber}: $f_I(s)\, d\theta \wedge \omega_I$.
\item \textbf{21 $d\chi$-fiber}: $f_I(s)\, d\chi \wedge \omega_I$.
\end{itemize}

T² isotropy $d\theta \leftrightarrow d\chi$ holds to $3 \times 10^{-7}$, and $35 + 21 + 21 = 77 = b_3$. The 35-block Gram matrix has condition number $7.66$ and eigenvalue range $[1.65, 12.62]$.

\textbf{Minimum associative cycle volume.} The 7 associative calibrated 3-cycles on $K_7$ (Fano-plane structure) have volumes from period integrals: $508.5,\; 362.3,\; 362.3,\; 219.9,\; 219.9,\; 219.9,\; 219.9$ (degeneracy pattern 1:2:4 reflecting $\mathbb{Z}_2$ Fano symmetry). The minimum satisfies an empirical scaling relation (no derivation currently available; presented as a conjectural numerical observation):
\begin{equation}
V_{\min} \approx \sqrt{\frac{\mathrm{Vol}(K_7)}{11}}, \qquad 11 = \frac{b_3}{n} = \frac{77}{7}.
\end{equation}
NK value: $219.90$; formula: $221.24$; deviation $0.6\%$. The factor $11 = b_3/7$ reflects the Fano-plane multiplicity, but no proof of this relation is available.

% ============================================
\section{Spectral Democracy (Numerical Verification)}\label{sec:democracy}
% ============================================

The spectral democracy theorem, proven analytically in [A] (Theorem~1.3), states that on product-type Ricci-flat $\Gtwo$ metrics with constant $g_{\theta\theta}$, the transverse 1-form Laplacian satisfies $\Delta_1(f\, d\theta) = (\Delta_0 f)\, d\theta$ exactly.

We verify this by comparing the first 60 eigenvalues of $\Delta_0$, $\Delta_1$ (transverse channel $f(s)\,d\theta$), and $\Delta_2$ (the corresponding 2-form channel). For each pair $(\Delta_p, \Delta_q)$, define the relative discrepancy:
\begin{equation}
\delta_{pq}(n) = \frac{|\lambda_n^{(p)} - \lambda_n^{(q)}|}{\lambda_n^{(p)}}.
\end{equation}

\begin{table}[ht]
\centering
\begin{tabular}{lcc}
\toprule
Comparison & $\max \delta$ ($n \leq 60$) & mean $\delta$ \\
\midrule
$\Delta_0$ vs $\Delta_1$ (transverse) & $8 \times 10^{-5}$ & $3 \times 10^{-5}$ \\
$\Delta_0$ vs $\Delta_2$ ($ds \wedge d\theta$ channel) & $9 \times 10^{-5}$ & $4 \times 10^{-5}$ \\
$\Delta_1$ vs $\Delta_2$ & $5 \times 10^{-5}$ & $2 \times 10^{-5}$ \\
\bottomrule
\end{tabular}
\caption{Spectral democracy discrepancy across Laplacian degrees.}
\label{tab:democracy}
\end{table}

All discrepancies are consistent with the fiber metric variation of $0.002\%$ ([A] Theorem~1.3 predicts exact agreement for a perfectly constant fiber). The spectral democracy is not generic: it would break down for metrics with significant $s$-dependent fiber geometry, and its observation at the $10^{-5}$ level provides an independent validation of the adiabatic decomposition.

The certified adiabatic bound ([A], Proposition~3.3) gives $\delta\lambda/\lambda \leq 0.78\%$; the observed $10^{-5}$ discrepancy is $80\times$ tighter, reflecting the T² isotropy ($3 \times 10^{-7}$) rather than the full adiabatic parameter ($2 \times 10^{-3}$).

% ============================================
\section{Discussion}\label{sec:discussion}
% ============================================

\subsection{Summary}

We have computed the spectral geometry of the NK-certified $\Gtwo$ metric of~[A]:
\begin{enumerate}
\item Scalar spectral gap $\lambda_1 = 0.12461 \pm 0.00016$ (Richardson) with Weyl law verification ($\alpha = 1.998$) and analytical formula $6\pi^2/475 \approx 0.12467$ (§\ref{sec:laplacian}).
\item Spectral confirmation of all Betti numbers with gap ratio $14{,}635$ for $b_2 = 21$ (§\ref{sec:harmonic}).
\item K3 fiber intersection form signature $(3, 18)$ (rank-21 restriction of K3 lattice) with SD/ASD gap $2{,}210$ (§\ref{sec:harmonic}).
\item Numerical confirmation of spectral democracy to $10^{-5}$ (§\ref{sec:democracy}).
\item Complete KK tower ($22{,}671$ levels) with Poisson statistics (§\ref{sec:laplacian}).
\item Conjectural scaling relation $V_{\min} \approx \sqrt{\mathrm{Vol}(K_7)/11}$ for minimum associative cycle volume, deviation $0.6\%$ (§\ref{sec:harmonic}).
\end{enumerate}

\subsection{Physical context}

The spectral data presented here is relevant for Kaluza--Klein reductions of supergravity theories on $\Gtwo$ manifolds. The scalar spectral gap determines the mass of the lightest KK mode; the harmonic forms determine the 4D field content; and the intersection form encodes gauge coupling structure.

A companion paper~[D] discusses the physical context of these spectral results in the $\Gtwo$ compactification setting.

\subsection{Limitations}

\textbf{Adiabatic approximation.} The spectral computations exploit the product-type structure of the certified metric ([A], Proposition~3.3). The adiabatic error is $\varepsilon_{\mathrm{ad}} = 2 \times 10^{-3}$ with certified eigenvalue perturbation $|\delta\lambda_n|/\lambda_n \leq 0.78\%$ for all eigenvalues. The spectral Betti confirmation (gap ratio $14{,}635$) is stable under this perturbation: a $0.78\%$ shift on each eigenvalue changes the gap ratio by at most $2 \times 0.78\% = 1.56\%$ (worst case, numerator and denominator moving in opposite directions), leaving it at four orders of magnitude. The spectral Betti identification $(b_2, b_3) = (21, 77)$ is a robust property of the torsion-free metric $g^*$ that survives any residual $O(10^{-3})$ uncertainty in the K3 fiber structural constants.

\textbf{Convergence.} The spectral gap is Richardson-extrapolated from 6 grids (Appendix~\ref{app:convergence}). The error bar $\pm 0.0001$ is a numerical estimate, not a rigorous bound.

\subsection{Open questions}

\begin{enumerate}
\item \textbf{Spectral characterization.} To what extent do the spectral data (gap, Weyl coefficient, SD/ASD ratio) characterize the underlying manifold among $\Gtwo$ 7-manifolds?
\item \textbf{Fiber-coupled corrections.} What is the spectral shift from explicit K3 mode coupling?
\item \textbf{Comparison with flow methods.} How does the spectrum compare with Laplacian flow metrics~\cite{LotayWei2019} on the same topological type?
\end{enumerate}

% ============================================
\section*{Appendix A. Convergence Analysis}\label{app:convergence}
\addcontentsline{toc}{section}{Appendix A. Convergence Analysis}
% ============================================

\begin{table}[ht]
\centering
\begin{tabular}{cccl}
\toprule
$N$ & $\lambda_1$ (Neumann) & $\lambda_1$ (Dirichlet) & $\Delta\lambda_1/\lambda_1$ \\
\midrule
200  & $0.12346$ & $0.12361$ & --- \\
400  & $0.12409$ & $0.12418$ & $0.51\%$ \\
600  & $0.12430$ & $0.12436$ & $0.17\%$ \\
800  & $0.12440$ & $0.12444$ & $0.08\%$ \\
1000 & $0.12446$ & $0.12449$ & $0.05\%$ \\
1200 & $0.12450$ & $0.12452$ & $0.03\%$ \\
\bottomrule
\end{tabular}
\caption{Richardson convergence of spectral gap across grid sizes and boundary conditions.}
\label{tab:convergence}
\end{table}

Extrapolated: $\lambda_1 = 0.12461 \pm 0.00016$.

% ============================================
\section*{Acknowledgements}
% ============================================

AI tools (Claude, Anthropic) were used for computational assistance. All scientific content, analysis, and conclusions are the sole responsibility of the author.

% ============================================
\section*{Author's note on AI collaboration}
% ============================================

This framework was developed through sustained collaboration between the author and several AI systems, primarily Claude (Anthropic), with contributions from GPT (OpenAI), Gemini (Google), and Aristotle (Harmonic) for specific mathematical insights. Architectural decisions and many key derivations emerged from iterative dialogue sessions. This collaboration follows a transparent crediting approach for AI-assisted mathematical research.

% ============================================
\section*{Disclosure statement}
% ============================================

The author reports there are no competing interests to declare.

% ============================================
\section*{Author's related works}
% ============================================

The following papers by the author are referenced above or provide directly related context. They are listed separately from the main bibliography to keep the citation record of this paper independent from the author's own prior output.

\begin{itemize}
\item[{[A]}] de La Fourni\`ere, B. \emph{A Certified Torsion-Free $\Gtwo$ Structure on a TCS Neck Model via Computer-Assisted Proof}. Zenodo, 2026. DOI: \href{https://doi.org/10.5281/zenodo.18860358}{10.5281/zenodo.18860358}.
\item[{[C]}] de La Fourni\`ere, B. \emph{Newton--Kantorovich diagnostics on a Donaldson K3 metric: two $\beta$ estimates, machine-checked inequalities}. Zenodo, 2026. DOI: \href{https://doi.org/10.5281/zenodo.19708916}{10.5281/zenodo.19708916}.
\item[{[D]}] de La Fourni\`ere, B. \emph{GIFT v3.4: geometric and physical context}. Companion paper, in preparation (2026).
\end{itemize}

% ============================================
\bibliographystyle{plain}
\bibliography{g2_spectral}
% ============================================

\end{document}
